\section{The integral adjoint argument}
\label{sec:integral-jump}

We prove Proposition~\ref{prop:integral-jump} using two uniform
estimates for forms.  The estimates are reduced in
Section~\ref{sec:tower} to Theorem~\ref{thm:bounded-fibre}, whose proof
occupies Sections~\ref{sec:bounded-presentations}--\ref{sec:incidence}.
All fields in these sections are finitely generated over $\C$; the
varieties in Proposition~\ref{prop:integral-jump} are transported to
$\C$ by the stipulated field isomorphisms.

A form of index $q$ on a function field is a nonzero rational section
of the $q$th tensor power of its absolute top differential line.
Its boundary on a normal model is minus its divisor divided by $q$.
On a model over a smaller function field, we divide by the $q$th
power of a base volume form to interpret this boundary relatively.
Rational convex combinations of form discrepancies are realized by
taking tensor powers and multiplying forms.  Indices of these
products need not be bounded unless a bound is expressly asserted.
A projective variety is of \emph{Fano type} if it admits an effective
klt boundary whose negative adjoint is ample.

\begin{theorem}[Lifting and Mixing]
\label{thm:lifting-mixing}
Fix positive integers $d,r$ and a positive real number $C$.
The following assertions hold.
\begin{enumerate}[label=\textup{(\arabic*)}]
\item\label{ij:lifting-item}
Let $P/Q$ be a regular extension with
$\operatorname{trdeg}_{\C}P\le d$, having a projective Fano type
generic model.  Let $\sigma,\psi$ be forms on $P$ whose boundaries on
this model are effective, with $\sigma$ of index at most $r$.
Suppose, on every nonzero divisorial valuation of $P$, that
\[
 A_\psi>0,\qquad 0\le A_\sigma\le C A_\psi.
\]
For every $\lambda>0$ there is
$\delta=\delta(d,r,C,\lambda)>0$ with the following property.
If $u$ is a nonzero integral-valued divisorial valuation of $Q$ and
\[
 \inf_{w|_Q=u}A_\psi(w)<\delta,
\]
then there is an integral-valued divisorial valuation $v$ of $P$
whose restriction is a positive multiple of $u$ and for which
$A_\psi(v)<\lambda$.  In relative dimension zero the regular
extension is $P=Q$, and the same assertion holds.
\item\label{ij:mixing-item}
Let $P$ have a projective Fano type model over $\C$ of dimension at
most $d$.  Let $\sigma,\psi,\chi$ be forms with effective boundaries
on that model, with $\sigma$ of index at most $r$, and suppose
\[
 A_\psi>0,\qquad 0\le A_\sigma\le C A_\psi.
\]
Suppose in addition that, for some $c>0$,
\[
 A_\chi(w)\ge c A_\psi(w)
 \quad\text{whenever }w\ne0\text{ and }A_\sigma(w)=0.
\]
There is a rational number $a\in(0,1)$ depending only on
$d,r,C,c$ such that
\[
 (1-a)A_\psi(w)+aA_\chi(w)>0
 \quad\text{for every nonzero divisorial }w.
\]
\end{enumerate}
All infima and homogeneous inequalities allow positive rational
scalings of prime divisorial valuations.  An integral-valued
valuation is not required to be primitive.
\end{theorem}

\subsection{The integral barrier}

\begin{proof}[Proof of Proposition~\ref{prop:integral-jump}, assuming
Theorem~\ref{thm:lifting-mixing}]
Write $A=A_{\eta_F}$ and $l_w=w(L)$.  The dimension of $F$ and the
clipping constant in Proposition~\ref{prop:integral-jump} are fixed
along the sequence.  We begin with a rational $t>0$, which will be
chosen only after obtaining constants independent of $t$.
Use the canonical divisor $K_F=\operatorname{div}(\eta_F)$ and put
\begin{equation}
\label{ij:adjoint-divisors}
 \Delta=\lceil tL\rceil-tL,
 \qquad E=K_F+\lceil tL\rceil.
\end{equation}
The boundary $\Delta$ has coefficients in $[0,1)$ and SNC support,
so $(F,\Delta)$ is klt.  Since $\eta_F$ is an lc volume form, every
coefficient of $K_F$ is an integer at least $-1$.  A coefficient
$-1$ gives $A(P)=0$ and hence $l_P>0$.  Thus $E$ is an effective
integral divisor.

Choose a sufficiently divisible $q$ and a general
$H_q\in|qtL|$, and set $H=H_q/q$.  Semiampleness permits $H_q$ to
be chosen smooth and transverse to the fixed SNC support.  Taking
$q\ge2$, the pair $(F,\Delta+2H)$ is sub-lc.  Consequently
\begin{equation}
\label{ij:general-member-bound}
 w(H)\le\tfrac12 A_{F,\Delta}(w)
 \quad\text{for every divisorial }w.
\end{equation}
Here and below discrepancies carrying a pair subscript are ordinary
pair discrepancies.  Define forms $\sigma,\psi$ by
\begin{equation}
\label{ij:two-forms}
 A_\sigma=A,
 \qquad A_\psi(w)=A(w)+t l_w-w(H).
\end{equation}
The first form is $\eta_F$, of index one; the second exists because
$tL-H\sim_{\Q}0$.  From
\[
 A_{F,\Delta}(w)=A(w)+t l_w-w(E)
\]
and \eqref{ij:general-member-bound}, we obtain
\begin{equation}
\label{ij:positive-anchor}
 A_\psi(w)\ge\tfrac12\bigl(A(w)+t l_w\bigr)>0,
 \qquad A_\sigma(w)\le2A_\psi(w).
\end{equation}
Strict positivity uses the hypothesis $l_w>0$ at every nonzero
divisorial valuation with $A(w)=0$.

The pair $(F,\Delta+H)$ is klt, its boundary is big, and its adjoint
is rationally linearly equivalent to $E\ge0$.  The good minimal
model theorem of \cite{BCHM} gives a $\Q$-factorial good minimal
model $Y$.  Let $h:Y\to S$ be its semiample contraction, with
connected fibres.  Subscripts denote birational transforms.
The rational principal divisor $tL-H$ is the same on every model,
so the ordinary MMP discrepancy inequality gives, on a common
resolution,
\begin{equation}
\label{ij:semiample-polarization}
 E_Y=K_Y+\Delta_Y+tL_Y=h^*N,
 \qquad w(E)\ge w(E_Y)=w(h^*N).
\end{equation}
Here $N$ is an effective ample $\Q$-Cartier divisor when
$\dim S>0$.  To obtain equality of divisors, take a globally
generated multiple of the effective semiample divisor $E_Y$.
Its defining section descends to the ample system on $S$ because
$h_*\cO_Y=\cO_S$, and its zero divisor defines that multiple of
$N$.  If $S$ is a point, the effective divisor $E_Y$ is rationally
linearly trivial and is zero; we then put $N=0$.
The effective pair $(Y,\Delta_Y+H_Y)$ is klt, and
\begin{equation}
\label{ij:ordinary-discrepancy}
 A_{Y,\Delta_Y+H_Y}(w)=A_\psi(w)-w(h^*N)>0.
\end{equation}

We claim that every integral-valued divisorial valuation satisfies
\begin{equation}
\label{ij:integral-barrier}
 w(h^*N)>0\quad\Longrightarrow\quad A_\sigma(w)+t l_w\ge1.
\end{equation}
If the conclusion fails, the nonnegative integer $A_\sigma(w)$ is
zero.  Let $D$ be the reduced SNC union of the supports of $K_F$
and $L$.  If $k_P$ is the coefficient of $K_F$ at a component of
$D$, then
\[
 0=A_\sigma(w)
   =A_{F,D}(w)+\sum_{P\subset D}(1+k_P)w(P).
\]
All summands are nonnegative.  Thus every component seen by $w$
has $k_P=-1$ and $l_P>0$.  Since $w(P)$ is a positive integer,
\[
 0<t l_P\le t l_w<1,
 \qquad \operatorname{coeff}_P(E)=-1+\lceil t l_P\rceil=0.
\]
The support of $E$ is contained in $D$, so $w(E)=0$, contradicting
\eqref{ij:semiample-polarization}.  This proves
\eqref{ij:integral-barrier}; its unit threshold is exactly where
the index-one hypothesis on $\eta_F$ is used.

\subsection{A uniformly positive base polarization}

Suppose $\dim S>0$, and put $P=\C(Y)$ and $Q=\C(S)$.
Connectedness makes $P/Q$ regular.  On the generic fibre of $h$,
the boundaries of $\sigma$ and $\psi$ are respectively
\[
 \Delta_Y+tL_Y,
 \qquad \Delta_Y+H_Y,
\]
restricted to that fibre, because $h^*N$ restricts to zero.
They are effective.  The second gives a klt log Calabi--Yau pair
with big boundary: $H_Y$ is big by birational pushforward and its
restriction to the generic fibre is big.  A klt log Calabi--Yau
pair with big boundary is of Fano type.  Explicitly, write its
boundary as a rational linear equivalence $B\sim_{\Q}A_0+D_0$
with $A_0$ ample and $D_0\ge0$.  For sufficiently small $e>0$,
$(1-e)B+eD_0$ is effective klt and its negative adjoint is
rationally linearly equivalent to $eA_0$.
If the generic fibre has dimension zero, regularity gives $P=Q$.

Apply Theorem~\ref{thm:lifting-mixing}(1) with reference index one,
comparison constant $2$, and, for example, $\lambda=1/4$.
By \eqref{ij:positive-anchor} and \eqref{ij:integral-barrier}, an
integral-valued valuation lying over a positive multiple of a
valuation $u$ with $u(N)>0$ has $A_\psi\ge1/2$.
Thus there is $\delta>0$, depending only on the fixed dimension
and independent of $t$, such that
\begin{equation}
\label{ij:base-lower-bound}
 \bar A_\psi(u):=\inf_{w|_Q=u}A_\psi(w)\ge\delta
 \quad\text{if }u\text{ is integral-valued and }u(N)>0.
\end{equation}
In relative dimension zero this follows directly from the barrier;
decreasing $\delta$ accommodates that case as well.

Apply the lc-trivial canonical bundle formula with b-semiample
moduli part, \cite[Theorems~1.5 and~6.28]{BFMT}, to
$h:(Y,\Delta_Y+H_Y)\to S$.
The morphism is projective and connected; the boundary is effective
and klt, and
\[
 K_Y+\Delta_Y+H_Y\sim_{\Q}h^*N.
\]
The discrepancy rank condition is one: on a generic resolution,
the rounded discrepancy divisor has no negative coefficients and
its positive coefficients occur only on exceptional divisors.
Its section space is therefore the constants on the connected
generic fibre.  These verify generic effectivity, generic
sub-log-canonicity, rank one, and the rational adjoint pullback
required by the cited theorem.

Its generalized base discrepancy is
\begin{equation}
\label{ij:base-generalized-discrepancy}
 \beta(u)=\bar A_\psi(u)-u(N)>0.
\end{equation}
Indeed, after putting $u$ on a smooth base model, the lc threshold
over its DVR is computed by resolving the fibre divisor together
with the boundary.  It is the minimum of the finitely many
discrepancy-to-multiplicity ratios of vertical components; the
horizontal directions are already klt.  Thus the infimum in
\eqref{ij:base-generalized-discrepancy} is attained and is positive.
On the original $S$ the discriminant is effective: over the generic
point of a prime of $S$ there is a prime of $Y$ of positive integral
fibre multiplicity, and its ordinary discrepancy is at most one.

On a high smooth model where the moduli divisor is semiample and
the crepant discriminant has SNC support, replace the moduli
divisor by $H_0/q_0$, where $H_0$ is a general member of its
$q_0$th multiple and $q_0\ge2$ is sufficiently divisible.
The subpair obtained by adding $H_0$ with coefficient one is
sub-lc, so $u(H_0)\le\beta(u)$ for all $u$.
After pushforward this gives an effective ordinary klt boundary
$B_S$ with
\begin{equation}
\label{ij:ordinary-base}
 K_S+B_S\sim_{\Q}N,
 \qquad \tfrac12\beta(u)\le A_{S,B_S}(u)\le\beta(u).
\end{equation}
Both the discriminant trace and the pushed-down general member
are effective.  On every model the replacement changes the adjoint
by the same rational principal divisor.  In particular
$K_S+B_S$ is $\Q$-Cartier and its ordinary crepant pullback is the
boundary just constructed.  This argument does not require
$K_S$ itself to be $\Q$-Cartier.

First take a small $\Q$-factorialization
$S_1\to S$ of the ordinary klt pair $(S,B_S)$.
Its crepant boundary is effective and pair discrepancies are
unchanged.  Now $K_{S_1}$ is $\Q$-Cartier and $S_1$ is klt, so
underlying-variety discrepancies are defined.  Fix
\[
 0<\epsilon'<\min\{1,\delta/4\},
\]
and extract to a $\Q$-factorial model $S'\to S_1$ precisely the
prime divisors whose underlying log discrepancy is less than
$\epsilon'$.  This is a finite set: on a fixed log resolution,
write the discrepancy as the nonnegative integral discrepancy of
the full reduced SNC support plus a positive linear combination
of its integral component orders.  A discrepancy below one forces
the integral term to vanish.  Such valuations are the monomial
lc places of the reduced SNC pair, and the positive coefficients
bound all their integral weights.  There are finitely many strata
and finitely many possible weight vectors.
The klt extraction furnished by \cite{BCHM} applies since each
desired divisor has pair discrepancy at most its underlying
discrepancy, hence less than one.  An explicit extraction
construction is given in Lemma~\ref{tower:geometric-step} below.

For this extraction $\pi:S'\to S_1$ one has
\[
 K_{S'}+B_{\rm exc}=\pi^*K_{S_1},\qquad B_{\rm exc}\ge0.
\]
Thus all remaining underlying discrepancies improve; the extracted
primes themselves have underlying discrepancy one on $S'$.
It follows that $S'$ is $\epsilon'$-lc.
The crepant transform $B_{S'}$ of the ordinary pair boundary is
effective.  If $N'$ denotes the pullback of $N$, then
\begin{equation}
\label{ij:birkar-hypotheses}
 N'\ge0,\qquad N'\text{ is nef and big},\qquad
 N'-K_{S'}\sim_{\Q}B_{S'}\ge0.
\end{equation}

Every positive coefficient of $N'$ is at least $\delta/2$.
For an extracted primitive valuation $u$ with $u(N)>0$,
\eqref{ij:ordinary-base} gives
\[
 \beta(u)\le2A_{S,B_S}(u)
 \le2A_{S_1}(u)<2\epsilon'.
\]
Together with \eqref{ij:base-lower-bound} this yields
$u(N)\ge\delta-2\epsilon'>\delta/2$.
For a prime already on $S_1$, use its corresponding prime on $S$.
Choose a prime of $Y$ dominating it, and let $k\ge1$ be its fibre
multiplicity.  The coefficient $k u(N)$ is a positive integer,
because $E_Y$ is the birational pushforward of the integral
divisor $E$.  Effectivity of $(Y,\Delta_Y+H_Y)$ gives
\[
 \beta(u)\le\frac1k\le u(N),
 \qquad
 \delta\le\bar A_\psi(u)=u(N)+\beta(u)\le2u(N).
\]
This proves the assertion for all primes of $S'$.

Birkar's effective birationality theorem
\cite[Theorem~4.2]{BirkarPolarised} now applies.
Its underlying variety is $\epsilon'$-lc; its polarization $N'$
is nef and big; $N'-K_{S'}$ is pseudo-effective; and in its required
decomposition into an integral pseudo-effective summand and an
effective summand we take $0+N'$.  The positive coefficients of
the latter are at least $\delta/2$.  We obtain a positive integer
$m$, taking a common multiple for the finitely many possible
positive dimensions of $S$, such that $|\lfloor mN'\rfloor|$ is birational and in
particular nonconstant.  Crucially, $m$ depends on neither $t$ nor
the denominators of the coefficients of $N'$.

\subsection{The strict ceiling inequality and final mixture}

We now fix a rational $t>0$ with $mt<1/20$.
Pass far enough along the sequence that, writing $\epsilon$ for
the prime estimate in Proposition~\ref{prop:integral-jump},
\begin{equation}
\label{ij:strict-choice}
 c_0:=m(t+\epsilon)<1/10.
\end{equation}
We may enlarge $\epsilon$ slightly to make it rational while
retaining this inequality.
The section $1$ belongs to $|\lfloor mN'\rfloor|$, and
nonconstancy gives another section represented by a nonconstant
rational function $g$ on $S'$.  Push its effective divisor
$\operatorname{div}(g)+mN'$ to $S$, pull back by $h$, and use
\eqref{ij:semiample-polarization}.  This gives
\begin{equation}
\label{ij:g-bound-on-F}
 \operatorname{div}(g)+mE\ge0
 \quad\text{on }F.
\end{equation}

Let $L_0$ be the effective big semiample $\Q$-Cartier divisor on
$F_0$ whose pullback is $L$.  At a prime of $F_0$ with $l_P=0$
we have $A(P)=1$, so its coefficient in the pushforward of $E$
is zero.  If $l_P>0$, its coefficient is
\[
 A(P)-1+\lceil t l_P\rceil
 \le A(P)+t l_P\le(t+\epsilon)l_P.
\]
Pushing \eqref{ij:g-bound-on-F} to $F_0$ therefore gives the
effective $\Q$-Cartier bound
\[
 \operatorname{div}(g)+c_0L_0\ge0.
\]
Its pullback controls every prime of $F$, including the exceptional
ones.  Set $b=1/10$.  If $l_P>0$ and
$n=\ord_P(g)\in\Z$, then
\[
 n\ge-c_0l_P>-b l_P
 \quad\Longrightarrow\quad
 n\ge1-\lceil b l_P\rceil.
\]
Since $K_{F,P}\ge-1$, this proves
$\operatorname{div}(g)+K_F+\lceil bL\rceil\ge0$ at such a prime.
If $l_P=0$, the pullback bound gives $n\ge0$, while
$K_{F,P}\ge0$: otherwise $K_{F,P}=-1$ would be an lc place with
$l_P=0$.  The same inequalities with $g=1$ hold as well.
Thus $1$ and the nonconstant $g$ are two independent sections of
$K_F+\lceil bL\rceil$, contrary to the assumed dimension one.
The positive-dimensional case is impossible, and for this fixed
$t$ we have $N=0$ sufficiently far along the sequence.

Now $(Y,\Delta_Y+H_Y)$ is a klt log Calabi--Yau pair with big
boundary, so $Y$ itself is of Fano type.  Let $J\ge0$ with
$J\sim_{\Q}L$ and define a form $\chi_0$ by
\[
 A_{\chi_0}(w)=A(w)+t l_w-tw(J).
\]
Its boundary on $Y$ is the effective divisor $\Delta_Y+tJ_Y$.
Write $C_0$ for the fixed clipping constant on the set $A=0$.
There, \eqref{ij:positive-anchor} and the hypothesis on $J$ give
\[
 A_{\chi_0}(w)\ge-(C_0-1)_+t l_w
 \ge-2(C_0-1)_+A_\psi(w),
\]
where $(x)_+=\max\{x,0\}$.  Choose a rational $\theta\in(0,1)$,
depending only on $C_0$, such that
\[
 1-\theta\bigl(1+2(C_0-1)_+\bigr)\ge\tfrac12.
\]
The form with discrepancy
$A_\chi=(1-\theta)A_\psi+\theta A_{\chi_0}$ has effective
boundary on $Y$ and satisfies $A_\chi\ge A_\psi/2$ on $A_\sigma=0$.
Theorem~\ref{thm:lifting-mixing}(2), with $r=1$, $C=2$ and
$c=1/2$, gives a uniform rational $a\in(0,1)$ such that, putting
$a_0=a\theta>0$,
\begin{align*}
 0
 &<(1-a)A_\psi(w)+aA_\chi(w)\\
 &=A(w)+t l_w-(1-a_0)w(H)-a_0t w(J).
\end{align*}
Since $H$ is effective, we conclude
\[
 w(J)\le\frac{A(w)+t l_w}{a_0t}
 \le\frac{\max\{1,t\}}{a_0t}\bigl(A(w)+l_w\bigr).
\]
The final constant is uniform because $t$ was fixed once and for
all.  This is the asserted integral jump estimate.
\end{proof}
